\documentclass[10pt,a4paper]{amsart}
\usepackage[margin=19mm]{geometry}
\usepackage{amsmath,amssymb,mathtools}
\usepackage[scr=rsfs,cal=euler]{mathalfa}
\usepackage{xfrac}
\usepackage[unicode,hidelinks,bookmarks=false]{hyperref}
\newcommand{\ie}{i.e.\ }
\newcommand{\cf}{cf.\ }
\newcommand{\resp}{respectively\ }
\newcommand{\Subsection}{\subsection}
\providecommand{\nobreakdash}{\nobreak\hbox{-}}
\setlength{\emergencystretch}{2em}
\multlinegap0pt
\usepackage{amssymb}
\usepackage{xcolor}
\usepackage{enumerate}
\usepackage{tikz}
\newtheorem{Assumption}{Assumption}
\newtheorem{theorem}{Theorem}
\title{On Lagrange multipliers of constrained optimization in Hilbert spaces}
\author{Zhiyu Tan}
\date{Source: arXiv:2306.03261v2 (CC BY 4.0)}
\begin{document}
\maketitle

\noindent\emph{Source and licence.} On Lagrange multipliers of constrained optimization in Hilbert spaces. Version arXiv:2306.03261v2 (CC BY 4.0), distributed by arXiv under the Creative Commons Attribution 4.0 International licence, http://creativecommons.org/licenses/by/4.0/, as recorded in the arXiv licence metadata for this exact version and shown on the abstract page https://arxiv.org/abs/2306.03261v2. Full licence notice: \texttt{licenses/arxiv-2306.03261-CC-BY-4.0.txt}.\par\smallskip

\noindent\textit{Editorial scope.} Counted unit: the article's theorem thm:pLM\_exist (source0.tex lines 550--569). The article's complete original proof of it is delivered with it (source0.tex lines 571--610, one \texttt{proof} environment of 40 lines). No mathematics is written, repaired or re-derived: the statement and the proof are the article's own lines, in the article's own notation, and no retained line is substituted or re-flowed.  Retained uncounted as the source-local context the proof needs: lines 539--541.  Nothing was omitted from any retained span, so the package carries no omission record.  No retained line contains a citation, so the wrapper carries no bibliography and no bibliography entry.  No retained line contains a figure environment, an image-inclusion command or drawing code, so no image is delivered and the package declares no asset dependency.  Editorial formatting only, in addition: an AMS \texttt{amsart} preamble that restates the article's own declarations and macros used by the retained lines, namely the environments \texttt{Assumption}, \texttt{theorem} on separate counters, together with the standard packages those lines need.  The retained environments print in this document as Assumption 1, Theorem 1 in order of appearance, whereas the article prints the same items with the numbering of its own full document.  The complete native source of the article as distributed in the arXiv e-print for this exact version is bundled unchanged as \texttt{sources/142535/source0.tex}.
\begin{Assumption}\label{assum:e_elm}
The essential Lagrange multiplier $\lambda^*$ exists at $u^*$ and $\lambda^* \neq 0$.
\end{Assumption}

\begin{theorem}\label{thm:pLM_exist}
Suppose that Assumption \ref{assum:e_elm} holds. The proper Lagrange multiplier exists at $u^*$ if and only if there exist  $\tilde{\lambda}\in \mathcal{N}(\zeta^*,K)$ and $\bar{\zeta}_0\in \overline{R(S)}$ such that
\begin{equation}\label{eq:LM_iff1}
 \mathrm{Ker}(\tilde{\lambda})\cap \overline{R(S)}=  \mathrm{Ker}(\lambda^*)\cap \overline{R(S)}\quad\quad (\mbox{Compatibility\ Condition})
 \end{equation}
 or equivalently
\begin{equation}\nonumber
\mathrm{Span}\{\tilde{\lambda}\} + \mathrm{Ker}(S^*) = \mathrm{Span}\{\lambda^*\} + \mathrm{Ker}(S^*)
 \end{equation}
and 
 \begin{equation}\label{eq:LM_iff2}
 \left\{
 \begin{aligned}
 & \langle\tilde{\lambda}, \bar{\zeta}_0 \rangle > 0,\\
 & \langle \lambda^*, \bar{\zeta}_0 \rangle > 0. 
 \end{aligned}
 \right.
 \quad\quad  (\mbox{Consistency\ Condition})
 \end{equation}
\end{theorem}

\begin{proof}
A direct calculation of the polar sets on both sides gives the equivalence.
%$$
%\mathrm{Ker}(\tilde{\lambda})\cap \overline{R(S)} = \mathrm{Ker}(\lambda^*)\cap \overline{R(S)} \ \Longleftrightarrow\ \mathrm{Span}\{\tilde{\lambda}\} + \mathrm{Ker}(S^*) = \mathrm{Span}\{\lambda^*\} + \mathrm{Ker}(S^*).
%$$
%\begin{equation}\nonumber
%\begin{aligned}
%&\quad\quad\quad\quad \mathrm{Ker}(\tilde{\lambda})\cap \overline{R(S)} = \mathrm{Ker}(\lambda^*)\cap \overline{R(S)}\\
%&\ \ \Longleftrightarrow\ \  [\mathrm{Ker}(\tilde{\lambda})\cap \overline{R(S)}]^\perp = [\mathrm{Ker}(\lambda^*)\cap \overline{R(S)}]^\perp\\
%&\ \ \Longleftrightarrow\ \  \mathrm{cl}\{ [\mathrm{Ker}(\tilde{\lambda})]^\perp + \mathrm{Ker}(S^*)\} = \mathrm{cl}\{[\mathrm{Ker}(\lambda^*)]^\perp + \mathrm{Ker}(S^*)\}\\
%&\ \ \Longleftrightarrow\ \ [\mathrm{Ker}(\tilde{\lambda})]^\perp + \mathrm{Ker}(S^*) = [\mathrm{Ker}(\lambda^*)]^\perp + \mathrm{Ker}(S^*)\\
%&\ \ \Longleftrightarrow\ \ \mathrm{Span}\{\tilde{\lambda}\} + \mathrm{Ker}(S^*) = \mathrm{Span}\{\lambda^*\} + \mathrm{Ker}(S^*).
%\end{aligned}
%\end{equation}

If the proper Lagrange multiplier $\bar{\lambda}$ exists, we have $\bar{\lambda} \in \mathcal{N}(\zeta^*,K)$ and $\lambda^* = \bar{\lambda}|_{\overline{R(S)}}$, which follows 
$$ \mathrm{Ker}(\bar{\lambda})\cap \overline{R(S)} = \mathrm{Ker}(\lambda^*)\cap \overline{R(S)}.$$
Since $\lambda^*\neq 0$ on $\overline{R(S)}$, there exists $\zeta_0\in \overline{R(S)}$ such that
$$\langle \lambda^*, \zeta_0\rangle \neq 0.$$
We assume that
$\langle \lambda^*, \zeta_0\rangle > 0.$
Otherwise, we can choose $-\zeta_0$. Taking $\bar{\zeta}_0 = \zeta_0$, the condition (\ref{eq:LM_iff2}) holds for $\bar{\lambda}$. 
Hence, we can take $\tilde{\lambda} = \bar{\lambda}$. 

Now we prove the other direction. 
Let $\tilde{\lambda}$ be an element in $\mathcal{N}(\zeta^*,K)$ such that (\ref{eq:LM_iff1}) and (\ref{eq:LM_iff2}) hold. 

Let  
$$ \bar{\lambda} = t_0\tilde{\lambda},$$
where $t_0 = \langle \lambda^*, \bar{\zeta}_0\rangle/\langle \tilde{\lambda}, \bar{\zeta}_0\rangle$ and $\bar{\zeta}_0$ satisfies (\ref{eq:LM_iff2}). 

We will prove that $\bar{\lambda}$ is a proper Lagrange multiplier. Note that the  condition (\ref{eq:LM_iff2}) implies $t_0>0$, which further gives $\bar{\lambda}\in \mathcal{N}(\zeta^*,K)$. 

Now, by the definition of the proper Lagrange multiplier,  we only need to show that
$$\lambda^* = \bar{\lambda}|_{\overline{R(S)}}.$$

By  (\ref{eq:LM_iff2}) and (\ref{eq:LM_iff1}), we have $\overline{R(S)} = \mathrm{Ker}(\lambda^*)\cap \overline{R(S)} + \mathrm{Span}\{\bar{\zeta}_0\} = \mathrm{Ker}(\tilde{\lambda})\cap \overline{R(S)} + \mathrm{Span}\{\bar{\zeta}_0\}$. Therefore, for any $\zeta \in \overline{R(S)}$, there exist $s\in \mathbb{R}$ and $\zeta_0\in \mathrm{Ker}(\tilde{\lambda})\cap \overline{R(S)}$ such that $\zeta = \zeta_0 + s\bar{\zeta}_0$. It follows
$$\langle \bar{\lambda}, \zeta\rangle =  t_0 \langle \tilde{\lambda}, \zeta\rangle =  t_0 \langle \tilde{\lambda}, \zeta_0 + s\bar{\zeta}_0\rangle = t_0s\langle \tilde{\lambda}, \bar{\zeta}_0\rangle = s\langle \lambda^*, \bar{\zeta}_0\rangle =  \langle \lambda^*, \zeta_0 + s\bar{\zeta}_0\rangle = \langle \lambda^*, \zeta\rangle,$$
which implies $\lambda^* = \bar{\lambda}|_{\overline{R(S)}}.$
\end{proof}

\end{document}
